\section{Ход исследования}

\subsection{Окружение}



Все эксперименты выполнены на одной машине с фиксированной конфигурацией,
чтобы обеспечить воспроизводимость и сопоставимость результатов.
\par

\textbf{Аппаратное обеспечение.}
\begin{itemize}
    \item Процессор: Intel Core Ultra 7 258V (8 ядер).
    \item Оперативная память: 32\,ГБ.
    \item Операционная система: Linux Ubuntu 25.04.
\end{itemize}

\textbf{СУБД.} PostgreSQL версии 18.2
\par

\textbf{Бенчмарк.}
В данной работе в качестве $\mathcal{Q}$ используются два индустриальных бенчмарка:
Join Order Benchmark (JOB) \cite{HowGoodJournal}~--- 113 запросов
с 5--17 соединениями к базе IMDb ~--- и JOB-Complex \cite{JOBComplex}~--- 30 запросов
к той же базе с реалистичными свойствами: соединения по строковым атрибутам, не первыичные, внешние ключи, 
сложные фильтры с множественными значениями.
IMDb (Internet Movie Database) содержит информацию о фильмах,
актёрах, ключевых словах и связанных сущностях.
Запросы JOB содержат от 5 до 17 соединений, коррелированные фильтры
и разнообразные топологии графов соединений,
что делает бенчмарк стандартным инструментом для оценки качества
планировщиков соединений.
\par

Эти свойства приводят к значительным ошибкам оценки кардинальностей:
PostgreSQL не имеет совместных статистик по строковым столбцам разных таблиц,
и предположение независимости нарушается на порядки.
Авторы \cite{JOBComplex} показали, что PostgreSQL выбирает планы
в 11{,}1 раза медленнее оптимальных на JOB-Complex (против 2 раз на JOB).
Использование JOB-Complex позволяет оценить устойчивость предложенных подходов
в условиях сильно ошибочных оценок кардинальностей.
\par

Полученные в работе результаты относятся к конкретным $\mathcal{Q}$ (JOB и JOB-Complex)
и не претендуют на существование универсального $\mathcal{A}^*$,
минимизирующего $t_{\mathrm{total}}$ для произвольного класса запросов.
\par

\textbf{Горячие данные.}
Все данные были загружены в буферный кэш СУБД. Это исключает влияние дискового ввода-вывода на измерения.


\textbf{Конфигурация СУБД.}
Параметры PostgreSQL выбраны так, чтобы устранить влияние побочных эффектов
(дисковый ввод-вывод, нехватка памяти для операторов) на время исполнения запросов,
изолируя эффект от выбора порядка соединений:
\begin{itemize}
    \item \texttt{shared\_buffers = 12\,GB}~--- размер буферного кэша СУБД.
    Достаточен для размещения всей базы IMDb в оперативной памяти.

    \item \texttt{work\_mem = 8\,GB}~--- объём памяти, выделяемый каждой
    операции сортировки или хеширования.

    \item \texttt{effective\_cache\_size = 8\,GB}~--- подсказка оптимизатору
    об объёме доступного кэша операционной системы.

    \item \texttt{max\_parallel\_workers = 8},
    \texttt{max\_parallel\_workers\_per\_gather = 8}~---
    параметры параллелизма.
\end{itemize}

\textbf{Метрики экспериментальной оценки.}
Прямое сравнение алгоритмов по абсолютному значению
$\sum_{q \in \mathcal{Q}} t_{\mathrm{total}}(\mathcal{A}, q)$
чувствительно к нескольким запросам с большим $t_{\mathrm{exec}}$,
которые могут доминировать в сумме и маскировать поведение на остальных.
Поэтому в работе используются три согласованные метрики, каждая из которых
отражает свой аспект близости алгоритма $\mathcal{A}$ к решению задачи
$\mathcal{A}^* = \arg\min_{\mathcal{A}} \sum_{q \in \mathcal{Q}} t_{\mathrm{total}}(\mathcal{A}, q)$:
\begin{itemize}
    \item \textbf{Нормализованное суммарное время}~--- отношение
    $\sum_q t_{\mathrm{total}}(\mathcal{A}, q)$/
    $\sum_q t_{\mathrm{total}}(\mathcal{A}_{\mathrm{DP}}, q)$,
    где $\mathcal{A}_{\mathrm{DP}}$~--- стандартный DPsize PostgreSQL.
    Значение $< 1$ означает, что $\mathcal{A}$ ближе к $\mathcal{A}^*$ по целевой функции,
    чем базовый алгоритм.

    \item \textbf{Нормализованные отношения по компонентам}~--- те же отношения
    отдельно для $t_{\mathrm{plan}}$ и $t_{\mathrm{exec}}$.
    Позволяют различать улучшения за счёт ускорения планирования
    и за счёт выбора лучшего плана.

    \item \textbf{Число, где лучше}~--- количество запросов $q \in \mathcal{Q}$,
    на которых $t_{\mathrm{total}}(\mathcal{A}, q) < t_{\mathrm{total}}(\mathcal{A}_{\mathrm{DP}}, q)$.
    Характеризует стабильность алгоритма: суммарное отношение может быть
    малым из-за нескольких крупных выигрышей при общей нестабильности.
\end{itemize}

\subsection{Подход 1. Жадные алгоритмы с различными критериями сравнения}

\textbf{Мотивация.}
Классическая реализация использует оценочную стоимость в качестве критерия.
Однако, как показано в секции <<Оценки в модели стоимости>>,
при увеличении числа соединений ошибки оценок кардинальностей
каскадируются и стоимость теряет способность
различать планы по истинному качеству.
Возникает вопрос: можно ли подобрать критерий сравнения,
более усточйчивый к ошибкам оценок?
\par

Реализованы и протестированы три варианта GOO, различающиеся
только критерием выбора лучшей пары на каждом шаге.
\par

\textbf{Вариант 1: GOO по стоимости} ($\mathrm{GOO_{cost}}$).
Критерий~--- оценочная стоимость результата соединения:
\[
    (R_i^*,\, R_j^*) \;=\; \arg\min_{(R_i, R_j)} \mathrm{Cost}(R_i \Join R_j).
\]
Это классическая реализация GOO, в которой на каждом шаге выбирается
пара с минимальной стоимостью по модели СУБД.
\par

\textbf{Вариант 2: GOO по кардинальности} ($\mathrm{GOO_{card}}$).
Критерий~--- оценка объёма результата соединения:
\[
    (R_i^*,\, R_j^*) \;=\; \arg\min_{(R_i, R_j)} \bigl(\widehat{|R_i \Join R_j|} \cdot w(R_i \Join R_j)\bigr),
\]
где $\widehat{|R_i \Join R_j|}$~--- оценочная кардинальность, а
$w(R_i \Join R_j)$~--- средняя ширина строки результата (в байтах).
Произведение $\widehat{|R_i \Join R_j|} \cdot w$ оценивает объём данных
промежуточного результата и является более грубым, но более устойчивым
критерием, чем стоимость на подпланах большого размера.
\par

\textbf{Вариант 3: GOO с комбинированным критерием} ($\mathrm{GOO_{comb}}$).
Критерий адаптивно переключается в зависимости от размера аргументов
соединения.
Пусть $|S|$~--- число базовых отношений в подплане $S$.
Пара $(R_i, R_j)$ называется \emph{простой}, если $|R_i| \leqslant 2$ и $|R_j| \leqslant 2$.
Критерий выбора:
\[
    (R_i^*,\, R_j^*) \;=\; \arg\min_{(R_i, R_j)}
    \begin{cases}
        \mathrm{Cost}(R_i \Join R_j), & \text{если пара простая},\\[4pt]
        \widehat{|R_i \Join R_j|} \cdot w(R_i \Join R_j), & \text{иначе}.
    \end{cases}
\]
\par

\textbf{Результаты.}
Все три варианта протестированы на 113 запросах JOB.
Суммарные отношения к DPsize:

\begin{center}
\begin{tabular}{l c c c c}
 & $t_{\mathrm{plan}}$ & $t_{\mathrm{exec}}$ & $t_{\mathrm{total}}$ & \#лучше \\
\hline
$\mathrm{GOO_{cost}}$ & $0{,}85$ & $1{,}44$ & $1{,}22$ & 55/113 \\
$\mathrm{GOO_{card}}$ & $0{,}86$ & $1{,}62$ & $1{,}34$ & 64/113 \\
$\mathrm{GOO_{comb}}$ & $0{,}81$ & $1{,}19$ & $1{,}05$ & 80/113 \\
\end{tabular}
\end{center}

Все три варианта сокращают $t_{\mathrm{plan}}$,
но проигрывают по $t_{\mathrm{exec}}$.
$\mathrm{GOO_{comb}}$ - наибольший суммарный выигрыш.
\par

\textbf{Случаи, когда жадные подходы превосходят DP.}
При значительных ошибках оценки кардинальностей DP может выбрать план,
оптимальный по ошибочным оценкам, но неудачный по фактическому времени.
Жадный алгоритм перебирает меньше вариантов, но его выбор на каждом шаге
менее подвержен каскадному накоплению ошибок, поскольку опирается
на локальное сравнение пар, а не на глобальную стоимость всего дерева.
\par

\textit{Пример: запрос 6d}
($t_{\mathrm{total}}^{\mathrm{DP}} = 2272$ мс,
$t_{\mathrm{total}}^{\mathrm{GOO_{comb}}} = 212$ мс, улучшение в $10{,}7$ раза).
Запрос содержит фильтр по таблице $n$ (name):
\texttt{name LIKE '\%Downey\%Robert\%'}, оставляющий 2~строки.
\par

DP строит левосторонний план $k \to mk \to t \to ci \to n$,
в котором таблица $n$ стоит последней.
Стоимостная модель недооценивает селективность шаблонного предиката \texttt{LIKE}
и оценивает промежуточный результат $ci$ как небольшой.
В действительности до соединения с $n$ кардинальность достигает
$785\,477$ строк, полученных сканированием по индексу.
\par

$\mathrm{GOO_{comb}}$ на первом шаге соединяет $n$ (2~строки после фильтра)
с $ci$ через индексное сканирование (243~строки),
поскольку эта пара имеет минимальную оценку кардинальности среди всех пар.
Далее результат соединяется хеш-соединением с веткой $k \to mk \to t$ ($14\,165$~строк).
Промежуточный результат после первого шага~--- 243~строки вместо $785\,477$.
Время исполнения: $117$ мс вместо $2163$ мс.
\par

Данный пример иллюстрирует общий механизм: когда одна из таблиц имеет
высокоселективный фильтр, который стоимостная модель не в состоянии
точно оценить, жадный выбор по кардинальности помещает её в начало
цепочки соединений, тогда как DP откладывает её на конец.
\par

\textbf{Случаи деградации жадных подходов.}
Жадный алгоритм принимает на каждом шаге необратимое решение
на основании локальной информации.
Если выбранное соединение создаёт промежуточный результат,
не позволяющий эффективно фильтровать на последующих шагах,
деградация неизбежна.
\par

\textit{Пример: запрос 17c}
($t_{\mathrm{total}}^{\mathrm{DP}} = 3086$ мс,
$t_{\mathrm{total}}^{\mathrm{GOO_{comb}}} = 10\,112$ мс,
деградация в $3{,}3$ раза).
На ранних шагах $\mathrm{GOO_{comb}}$ соединяет $k \to mk \to t \to mc \to cn$:
каждая пара имеет малую оценку кардинальности, и на каждом шаге
выбор выглядит оптимальным.
Однако построенная ветвь не содержит таблицу $n$ (name), имеющую
высокоселективный фильтр \texttt{name LIKE '\%B\%'}.
Когда на последних шагах присоединяются $ci$ ($52{,}5$ строки на итерацию,
$148\,552$ итераций) и затем $n$ ($7\,796\,926$ итераций),
фильтрация по $n$ применяется поздно, и промежуточные результаты
оказываются на порядок больше, чем в плане DP.
\par

DP находит порядок, в котором $ci$ соединяется с $n$ рано
(ещё до $mc$ и $cn$), что позволяет фильтру по имени сократить
промежуточный результат до $250$~строк.
Жадный алгоритм не способен предвидеть этот эффект,
поскольку на момент принятия решения таблица $n$ не связана
предикатом с текущим промежуточным результатом.
\par

\textbf{Результаты на JOB-Complex.}
На 30 запросах JOB-Complex картина меняется радикально:

\begin{center}
\begin{tabular}{l c c c c}
 & $t_{\mathrm{plan}}$ & $t_{\mathrm{exec}}$ & $t_{\mathrm{total}}$ & \#лучше\\
\hline
$\mathrm{GOO_{cost}}$ & $0{,}81$ & $1{,}92$ & $1{,}85$ & 9/30 \\
$\mathrm{GOO_{card}}$ & $0{,}86$ & $0{,}06$ & $0{,}11$ & 17/30 \\
$\mathrm{GOO_{comb}}$ & $0{,}85$ & $0{,}13$ & $0{,}17$ & 16/30 \\
\end{tabular}
\end{center}

$\mathrm{GOO_{card}}$ показывает улучшение $t_{\mathrm{total}}$ в 9 раз
относительно DPsize~--- результат, диаметрально противоположный JOB,
где $\mathrm{GOO_{card}}$ проигрывал ($t_{\mathrm{total}} = 1{,}34$).
$\mathrm{GOO_{comb}}$ также улучшает $t_{\mathrm{total}}$ в $5{,}9$ раз.
Напротив, $\mathrm{GOO_{cost}}$ деградирует ($1{,}85$), подтверждая, что на JOB-Complex
стоимостная модель сильно ошибается.
\par

JOB-Complex содержит соединения по строковым столбцам без первичных/внешних ключей,
для которых ошибки оценки кардинальностей достигают $10^5$--$10^8$.
При таких ошибках DP выбирает планы со вложенными циклами, оптимальные по стоимости, 
но долгие по фактическому времени.
Например, запрос~12: DPsize выбирает вложенный цикл с $2{,}6 \cdot 10^9$ итераций
($t_{\mathrm{exec}} = 539$ с), тогда как $\mathrm{GOO_{card}}$ находит
хеш-соединение ($t_{\mathrm{exec}} = 640$ мс, улучшение в 840 раз).
$\mathrm{GOO_{card}}$, поскольку не использует
стоимостную модель и не может выбрать вложенный цикл на основании
ошибочно заниженной оценки кардинальности.
\par

\subsection{Подход 2. Итеративная декомпозиция графа соединений на топологии}



\textbf{Мотивация.}
Как отмечается в \cite{SizeSub, Adaptive}, время работы алгоритмов перечисления
подпланов существенно зависит от топологии графа соединений.
Для цепи из $n$ вершин DPccp перебирает $O(n^3)$ пар связных подграфов,
для звезды~--- $O(n \cdot 2^n)$, для полного графа~--- $O(3^n)$.
Возникает естественная идея: если граф запроса можно разложить на подграфы
с известной топологией, каждый можно спланировать специализированным алгоритмом
с меньшей сложностью, чем общий DP для всего графа.
\par

\textbf{Описание подхода.}
Граф соединений итеративно разбивается на топологические фрагменты
(плотные подграфы, звёзды, циклы, цепи)~--- в фиксированном порядке приоритетов.
Каждый фрагмент планируется специализированным алгоритмом: цепи и циклы~---
интервальным DP, звёзды~--- планированием лучей с жадной сборкой к центру,
плотные подграфы~--- DPsub. Результат планирования фрагмента становится 
вершиной нового графа; процесс повторяется до схождения к
одной вершине. Финальная сборка компонент выполняется жадным
алгоритмом $\mathrm{GOO_{comb}}$.
\par

\textbf{Результаты.}
Подход протестирован на всех 113 запросах JOB и 30 запросах JOB-Complex.

\begin{center}
\begin{tabular}{l c c c c}
Бенчмарк & $t_{\mathrm{plan}}$ & $t_{\mathrm{exec}}$ & $t_{\mathrm{total}}$ & \#лучше \\
\hline
JOB         & $0{,}83$ & $3{,}55$ & $2{,}55$ & 32/113 \\
JOB-Complex & $0{,}82$ & $3{,}12$ & $2{,}99$ & 12/30 \\
\end{tabular}
\end{center}

Экономия $t_{\mathrm{plan}}$ в 17\% достигается за счёт меньшего числа вызовов
процедуры построения путей в каждом фрагменте, но деградация
$t_{\mathrm{exec}}$ в $3{,}1$--$3{,}6$ раза её полностью перекрывает.
Из 113 запросов JOB лишь 32 показали улучшение $t_{\mathrm{total}}$;
81 запрос деградировал, причём 10 запросов~--- более чем в 5 раз.
\par

\textbf{Причина деградации: потеря параметризованных путей через центральные вершины.}
В нормализованных реляционных схемах существуют центральные таблицы,
к которым через внешние ключи присоединяются многие другие.
В схеме IMDb такой вершиной является таблица $t$ (title):
к ней присоединяются $ci$, $mc$, $mk$, $mi$, $mi\_idx$, $cc$.
При декомпозиции центральная попадает в один фрагмент, а соседние таблицы
из других фрагментов теряют возможность использовать индексное сканирование
по первичному ключу таблицы.
Вместо параметризованного индексного сканирования 
оптимизатор вынужден использовать хеш-соединение без индекса.
\par

\textbf{Пример: запрос 10a.}
Запрос соединяет 7 таблиц, включая центральную вершину $t$ и большие таблицы
$ci$, $mc$ (кардинальности ${\sim}36$~млн, ${\sim}2{,}5$~млн, ${\sim}2{,}6$~млн).
Стандартный DPsize строит левостороний план с промежуточными
кардинальностями $681 \to 4395 \to 2270 \to 76 \to 56 \to 52 \to 52$,
используя индексное сканирование $t$ по первичному ключу при каждом
присоединении ($t_{\mathrm{exec}} = 176$ мс).
\par

Алгоритм декомпозиции относит $t$ в один фрагмент с $mc$, а $ci$~---
в другой фрагмент с $rt$ и $chn$. После планирования фрагментов 
новая вершина $\{mc, t\}$ содержит $207\,410$ строк (хеш-соединение
двух полностью сканированных таблиц), а вершина $\{rt, ci, chn\}$~---
$25\,523$ строки. На этапе сборки эти две новые вершины соединяются
хеш-соединением по $t.\mathit{id} = ci.\mathit{movie\_id}$,
обрабатывая $207\,410 \times 25\,523$ кандидатов вместо 76 строк
параметризованного индексного сканирования в плане DPsize.
$t_{\mathrm{exec}}$ растёт с $176$ мс до $2{,}66$ секунд (деградация в $15{,}1$ раза).
Аналогичная потеря параметризованных путей через $t$ наблюдается в 78 из 81
деградировавших запросов.
\par

\textbf{Вывод.}
Подход с декомпозицией на топологии бесперспективен для схем с выраженными
центральными вершинами, характерными для нормализованных реляционных баз данных.
Корень проблемы~--- разрушение параметризованных индексных путей при
изоляции центральных вершин в один фрагмент.
Это наблюдение мотивирует \emph{безопасную} стратегию сокращения пространства
поиска, в которой центральные не могут быть изолированы от своих соседей.



\subsection{Подход 3. DPhyp (реализация с открытым исходным кодом)}



\textbf{Мотивация.}
Алгоритм DPhyp \cite{MoeNeu} достигает нижней теоретической границы перебора
среди DP-алгоритмов для произвольных графов соединений:
он перечисляет \emph{только} пары связных подграфов (csg-cmp-pairs),
тогда как DPsize перебирает все подмножества, проверяя связность апостериори.
Для разреженных графов (цепи, звёзды) это может дать экспоненциальное
сокращение числа рассматриваемых пар.
\par

Помимо теоретической границы перебора, DPhyp меняет \emph{порядок},
в котором кандидаты поступают в процедуру построения путей:
DPsize обрабатывает подпланы по возрастанию размера, а DPhyp~--- обходом
по соседям гиперграфа.
Поскольку результат планирования в PostgreSQL зависит от порядка перебора
(см. <<Архитектура планировщика PostgreSQL>>),
эксперимент позволяет одновременно проверить две гипотезы:
(1) приводит ли достижение нижней теоретической границы перебора
к улучшению $t_{\mathrm{plan}}$ и (2) может ли альтернативный порядок обхода
систематически улучшить $t_{\mathrm{exec}}$.
\par

\textbf{Описание.}
Использована реализация DPhyp с открытым исходным кодом,
интегрированная в PostgreSQL в виде расширения.
Алгоритм строит гиперграф соединений и перечисляет пары
связных подграфов $(S_1, S_2)$ обходом по соседям, гарантируя,
что каждая перечисленная пара связна и соединена хотя бы одним
предикатом. Для каждой пары вызывается процедура построения путей,
аналогичная стандартному DPsize.
\par

\textbf{Результаты.}
Суммарные отношения к стандартному DPsize:
\begin{itemize}
    \item $t_{\mathrm{plan}}$: $1{,}015$ (деградация 1{,}5\%),
    \item $t_{\mathrm{exec}}$: $1{,}131$ (деградация 13{,}1\%),
    \item $t_{\mathrm{total}}$: $1{,}088$ (деградация 8{,}8\%).
\end{itemize}
Алгоритм, достигающий нижней теоретической границы перебора,
показал результат \emph{хуже} стандартного DPsize по всем трём метрикам.
\par

\textbf{Результаты на JOB-Complex.}
На 30 запросах JOB-Complex DPhyp также уступает DPsize:
$t_{\mathrm{plan}} = 1{,}06$, $t_{\mathrm{exec}} = 1{,}10$,
$t_{\mathrm{total}} = 1{,}10$, \#лучше 8 из 30.
При этом на отдельных запросах с плохими планами DPsize
(запросы~5, 11, 12, 20, где DPsize выбирает планы с вложенными циклами с $10^9$ итераций)
DPhyp уступает им в $1{,}5$--2 раза, тогда как $\mathrm{GOO_{comb}}$
находит на порядок более быстрый план.
Это показывает, что изменённый порядок перечисления DPhyp \emph{не устраняет}
сильные ошибки стоимостной модели на JOB-Complex, а меняет их распределение.
\par

\textbf{Анализ причин.}
\par

\textit{1. Теоретическое преимущество не реализуется на практике.}
Свыше 90\% времени $t_{\mathrm{plan}}$ приходится на процедуру
построения путей, а не на само перечисление пар.
DPhyp и DPsize вызывают эту процедуру для одних и тех же допустимых пар,
различаясь лишь порядком вызовов.
\par

\textit{2. Изменённый порядок перечисления ухудшает $t_{\mathrm{exec}}$.}
Это главная причина деградации.
DPhyp перечисляет пары обходом по соседям гиперграфа,
а DPsize~--- в порядке возрастания размера подмножества.
Как показано в секции <<Архитектура планировщика PostgreSQL>>,
при ошибках оценок кардинальностей стоимости конкурирующих путей
попадают в $\varepsilon$-окрестность, и путь, поступивший
в процедуру инкрементального отбора раньше, получает структурное
преимущество.
DPsize на каждом уровне сначала обрабатывает левостороние планы, где на каждом уровне 
присоединяются единичны отношения, для которых стоимостная модель PostgreSQL наиболее
точна.
DPhyp не имеет этого свойства: порядок перечисления определяется
структурой гиперграфа и может отдавать приоритет ветвистым планам
с менее надёжными оценками.
\par

\textbf{Вывод.}
Эксперимент с DPhyp показывает, что замена алгоритма перечисления
при сохранении пространства поиска и процедуры построения путей
не даёт систематического улучшения $t_{\mathrm{total}}$.
Причина~--- механизм $\varepsilon$-зависимости от порядка перебора:
при ошибочных оценках кардинальностей стоимости конкурирующих путей
попадают в $\varepsilon$-окрестность, и изменение порядка перечисления
может как улучшить, так и ухудшить отдельные запросы, без устойчивого
преимущества. Более того, если новый порядок систематически смещает
приоритет в сторону планов с менее надёжными оценками, результат
становится хуже DPsize.
\par

Отсюда следуют два структурных ограничения для $t_{\mathrm{total}}$,
не устранимые заменой алгоритма перечисления:
\begin{itemize}
    \item $t_{\mathrm{exec}}$ сильно зависит от качества оценок кардинальностей
    и порядком поступления путей в процедуру инкрементального отбора,
    а не алгоритмом обнаружения пар.
    \item $t_{\mathrm{plan}}$ определяется числом вызовов процедуры
    построения путей, которое при неизменном пространстве поиска
    совпадает у DPsize и DPhyp с точностью до фильтрации несвязных пар.
\end{itemize}
Для улучшения $t_{\mathrm{total}}$ необходимо либо сократить пространство
поиска (фиксируя соединения с предсказуемыми оценками до основного DP),
либо целенаправленно упорядочить перебор так, чтобы приоритет
получали планы с надёжными оценками. Первое направление реализовано
в подходе со связыванием якорных участков (подход~5), второе~--- в модификации
DPhyp, описанной далее (подход~4).
\par

Следует подчеркнуть, что результат характеризует данную реализацию DPhyp,
а не алгоритм в целом. Полноценная реализация с корректным построением
обобщённых гиперрёбер и обнаружением конфликтов может показать иные результаты.
Тем не менее пункты 1 (теоретическое преимущество не реализуется при $n \leqslant 17$)
и 2 (зависимость от порядка перечисления) являются свойствами архитектуры
PostgreSQL, а не конкретной реализации DPhyp.


\subsection{Подход 4. Модифицированный DPhyp (DPhypMod)}



\textbf{Мотивация.}
Эксперимент с DPhyp (подход~5) выявил три конкретные причины деградации:
нарушение поуровневого инварианта PostgreSQL, невыгодный порядок
обработки кандидатов и избыточные вызовы процедуры построения путей
на промежуточных уровнях. Модифицированный DPhyp адресует каждую из этих
причин отдельной модификацией, сохраняя оригинальный алгоритм
перечисления csg-cmp-пар без изменений.
\par

\textbf{Модификация 1. Поуровневая материализация.}
Стандартный DPhyp строит отношения в порядке обхода гиперграфа,
произвольно чередуя подпланы разных уровней.
Однако архитектура планировщика PostgreSQL опирается на инвариант:
все отношения уровня $k$ должны быть полностью построены
и финализированы до начала построения отношений уровня $k+1$.
Нарушение этого инварианта приводит к неполному распространению
классов эквивалентности и пропуску путей.
\par

Модификация разделяет работу DPhyp на две фазы:
\begin{enumerate}
    \item \textbf{Перечисление.} Оригинальный алгоритм DPhyp перечисляет
    все csg-cmp-pairs, но не вызывает процедуру построения путей.
    Вместо этого для каждой пары $(S_1, S_2)$ записывается
    список кандидатов в соответствующий узел DP-таблицы
    и вычисляется облегчённая оценка кардинальности $\widehat{|S_1 \Join S_2|}$.
    \item \textbf{Материализация.} Узлы DP-таблицы группируются по уровню
    (числу базовых отношений) и обрабатываются строго последовательно:
    сначала все узлы уровня $2$, затем уровня $3$, и~т.\,д.
    Для каждого узла вызывается процедура построения путей по всем
    его кандидатам, после чего выполняется финализация.
    Это восстанавливает поуровневый инвариант PostgreSQL.
\end{enumerate}

\textbf{Модификация 2. Приоритет лесторонние планов.}
Как показано в секции <<Архитектура планировщика PostgreSQL>>,
стандартный DPsize на каждом уровне обрабатывает
левостороние подпланы $(|S_1| = s-1, |S_2| = 1)$ до ветвистых.
Это обеспечивает приоритет параметризованных путей
(вложенный цикл с индексным сканированием) в процедуре
инкрементального отбора.
В DPhyp порядок кандидатов определяется обходом гиперграфа
и не имеет этого свойства.
\par

Модификация вводит два уровня сортировки:
\begin{enumerate}
    \item \textbf{Между узлами одного уровня.}
    Узлы уровня $k$ группируются по значению $\mu$~---
    минимальному размеру дочернего подплана среди всех кандидатов узла.
    Узлы с $\mu = 1$ (левостороние) обрабатываются первыми,
    затем $\mu = 2$, и~т.\,д.
    \item \textbf{Внутри списка кандидатов одного узла.}
    Пары кандидатов сортируются по $\min(|S_1|, |S_2|)$
    по возрастанию, так что левосторонние планы узла
    поступают в процедуру построения путей раньше ветвистых.
\end{enumerate}
Обе сортировки воспроизводят порядок обработки стандартного
DPsize в контексте DPhyp-перечисления.
\par

\textbf{Модификация 3. Отсечение на промежуточных уровнях по кардинальности.}
Распределение числа узлов DP-таблицы по уровням имеет характерную
форму нормального распределения с максимумом на средних уровнях
($s \approx n/2$). Для запроса с $n = 11$ таблицами на уровнях $5$--$8$
сосредоточено ${\sim}86\%$ всех узлов. Основная вычислительная нагрузка
планирования приходится именно на эти уровни.
\par

На каждом уровне из области горба ($\lceil n/3 \rceil + 1 \leqslant s \leqslant n - \lfloor n/3 \rfloor$)
выполняется отсечение узлов с высокой оценкой кардинальности:
\begin{enumerate}
    \item Вычисляется $\min_{S}\ \widehat{|S|}$ среди всех узлов уровня $s$.
    \item Узлы с $\widehat{|S|} > \alpha \cdot \min_{S}\ \widehat{|S|}$
    отсекаются (в реализации $\alpha = 300$).
    \item Дополнительно отсекаются узлы с $\widehat{|S|} > C_{\mathrm{out}}^{\mathrm{GOO}}$,
    где $C_{\mathrm{out}}^{\mathrm{GOO}}$~--- сумма кардинальностей промежуточных результатов
    жадного плана, вычисленная облегчённым GOO по оценкам $\widehat{|S|}$
    из фазы перечисления.
    Если кардинальность одного промежуточного результата превышает
    суммарный объём всех промежуточных результатов жадного плана,
    такой подплан заведомо неоптимален.
    \item \textbf{Гарантия покрытия.} После отсечения проверяется,
    что каждое базовое отношение $R_i \in \mathcal{R}$ входит хотя бы
    в один выживший узел уровня $s$.
    Если отношение потеряно, восстанавливается узел с минимальной
    $\widehat{|S|}$ среди отсечённых, содержащий это отношение.
    Это предотвращает отсечение всех путей через критически важные
    базовые отношения.
\end{enumerate}

На крайних уровнях ($s \leqslant \lceil n/3 \rceil$ и $s > n - \lfloor n/3 \rfloor$)
отсечение не выполняется: число узлов мало, стоимостная модель
ещё точна, а разнообразие путей критически важно для качества
финального плана.



\textbf{Результаты.}
Суммарные отношения к стандартному DPsize:
\begin{itemize}
    \item $t_{\mathrm{plan}}$: $0{,}966$ (экономия 3{,}4\%),
    \item $t_{\mathrm{exec}}$: $0{,}967$ (улучшение 3{,}3\%),
    \item $t_{\mathrm{total}}$: $0{,}967$ (улучшение 3{,}3\%).
\end{itemize}
Из 113 запросов 98 показали улучшение $t_{\mathrm{total}}$,
что значительно лучше, чем исходный DPhyp (78 из 113).
\par

Абсолютные суммарные времена (по 113 запросам):
\begin{itemize}
    \item DPsize: $t_{\mathrm{plan}} = 52{,}3$ с,\;
    $t_{\mathrm{exec}} = 88{,}8$ с,\; $t_{\mathrm{total}} = 141{,}2$ с.
    \item DPhypMod: $t_{\mathrm{plan}} = 50{,}5$ с,\;
    $t_{\mathrm{exec}} = 85{,}9$ с,\; $t_{\mathrm{total}} = 136{,}5$ с.
    \item DPhyp: $t_{\mathrm{plan}} = 53{,}1$ с,\;
    $t_{\mathrm{exec}} = 100{,}5$ с,\; $t_{\mathrm{total}} = 153{,}6$ с.
\end{itemize}

\textbf{Анализ: почему улучшение ограничено 3{,}3\%.}
\par

\textit{1. Пространство поиска совпадает с DPsize.}
Все три модификации влияют на \emph{порядок обработки} и \emph{количество}
узлов DP-таблицы, но не расширяют и не сужают множество рассматриваемых
разбиений по сравнению с DPsize.
Поуровневая материализация и сортировка кандидатов восстанавливают
инварианты стандартного планировщика~--- но не превосходят их.
DPhypMod \emph{приближается} к DPsize по качеству планов,
а не улучшает его.
\par

\textit{2. Отсечение экономит мало.}
Запросы JOB содержат от 5 до 17 таблиц.
При таких размерах число узлов DP-таблицы составляет от десятков до нескольких
тысяч, а горб распределения~--- область, где отсечение действует~---
содержит от десятков до нескольких сотен узлов.
Отсечение по кардинальности с консервативным порогом $\alpha = 300$
и гарантией покрытия удаляет лишь единицы узлов на каждом уровне.
Экономия $t_{\mathrm{plan}}$ от сокращения вызовов процедуры построения путей
составляет 3{,}4\%, что соответствует удалению ${\sim}3\%$ узлов.
Более агрессивное отсечение ($\alpha < 100$) рискует удалить узлы,
через которые проходит оптимальный план, поскольку оценки кардинальностей
на промежуточных уровнях ненадёжны.
\par

\textit{3. Процедура построения путей остаётся узким местом.}
Свыше 90\% времени $t_{\mathrm{plan}}$ приходится на процедуру построения путей,
а не на перечисление или сортировку узлов.
Замена алгоритма перечисления и оптимизация порядка
обработки не могут существенно сократить число вызовов этой процедуры,
потому что число допустимых разбиений определяется топологией графа
соединений.
\par

Для существенного улучшения необходимо сократить пространство поиска~--- уменьшая число отношений,
подаваемых на вход DP, путём предварительного связывания подграфов
с предсказуемыми свойствами.
Этот вывод мотивирует переход к подходу со связыванием якорных участков.

\textbf{Результаты на JOB-Complex.}
На JOB-Complex: $t_{\mathrm{plan}} = 1{,}04$,
$t_{\mathrm{exec}} = 0{,}82$, $t_{\mathrm{total}} = 0{,}83$, \#лучше 10 из 30.
Улучшение $t_{\mathrm{exec}}$ на 18\% объясняется тем, что поуровневая
материализация и сортировка кандидатов частично восстанавливают
инварианты DPsize, но не устраняют значительные ошибки оценок
на строковых соединениях.
На запросах с умеренными ошибками (Q8, Q9, Q13, Q14) DPhypMod
показывает результаты, сопоставимые с DPsize.
На запросах с выраженной деградацией (Q5, Q11, Q12, Q20)
он не достигает уровня $\mathrm{GOO_{card}}$ и $\mathrm{GOO_{comb}}$,
поскольку отсечение по $\widehat{|S|}$ на промежуточных уровнях
неэффективно, когда сами оценки ошибочны на порядки.
\par

\subsection{Подход 5. Связывание якорных участков}



\textbf{Мотивация.}
Исходя из аннотации и выводов с прошлых подходов следует, 
что адаптивное планирование имеет смысл рассматривать как 
\emph{предобработка графу} перед основным DP:
сокращение числа вершин за счёт стягивания подграфов с надёжными
оценками кардинальностей, после которой стандартный DPsize
работает быстрее и устойчивее.
Переключение между стратегиями происходит по единственному критерию
$\#\mathrm{csg}$, естественным образом отражающему обещанную
в аннотации <<сложность графа>>: для простых графов применяется
прямой DPsize, для сложных~--- связывание с последующим DPsize.
\par

При этом связывание должно быть \emph{безопасной}: стягиваться
могут только те подграфы, для которых оценки кардинальностей
заведомо надёжны, а центральные вершины не могут быть изолированы
от своих соседей (это требование~--- прямое следствие провала
подхода~2).
\par

\textbf{Определения.}
Пусть задан граф соединений $G = (\mathcal{R}, \mathcal{E})$ с $n$ вершинами.
\begin{itemize}
    \item \textbf{Якорная вершина}~--- вершина $v \in \mathcal{R}$, удовлетворяющая
    двум условиям:
    \begin{enumerate}
        \item $|v| < R_{\max}$
        \item $|v| \cdot \alpha < \max_{u \in N(v)} |u|$
    \end{enumerate}
    Якорь~--- это малая таблица, имеющая хотя бы одного соседа,
    который в $\alpha$ раз больше. Соединение якоря с крупным соседом высокоселективно,
    а его кардинальность оценивается надёжно, поскольку на единичных отношениях 
    оценки более точные.

    \item \textbf{Якорный участок}~--- связная компонента подграфа, включающая в себя якорную вершину 
    и её соседей. Участок имеет ограниченную кардинальность:
    \begin{enumerate}
        \item $|v \cup N(v)| \leqslant L$, где v - якорная вершина
    \end{enumerate}
    Если якорные участки смежны, т.е существует условие соединения между двумя вершинами 
    из разных участков, то два участка объединяются в один.
\end{itemize}

\textbf{Алгоритм.}
Входом является список базовых отношений и граф соединений.
Алгоритм реализует каскад из трёх стратегий с убывающей точностью:
\par

\textit{Шаг 1. Построение графа и оценка сложности.}
Строится граф соединений $G$. Для каждой связной компоненты 
$C$ в $G$ вычисляется $\#\mathrm{csg}(C)$~---
число связных подграфов, определяющее объём DP-таблицы и сложность перебора.
Если $\#\mathrm{csg}(C) \leqslant \tau$,
граф передаётся стандартному DPsize без модификаций:
при малом $\#\mathrm{csg}$ точный перебор быстр.
\par

\textit{Шаг 2. Связывание якорных участков.}
Если $\#\mathrm{csg}(G) > \tau$, выполняется поиск и связание N итераций. Если на каком-то шаге условие 
перестаёт выполняться, то происходит досрочный выход из цикла итераций:
\begin{enumerate}
    \item Для каждой вершины проверяется якорное условие.
    Все якори и их соседи маркируются.
    \item Удаление небезопасных вершин. Если вершина не является якорем, но входит в участок и содержит строковое 
    соединение, то она удаляется из участка
    \item Маркированные вершины разбиваются
    на связные компоненты обходом в ширину по подграфу маркированных
    вершин. Каждая связная компонента образует один якорный участок.
    \item Каждый участок планируется \emph{стандартным DPsize}.
    \item Если кардинальность результата планирования участка
    превышает порог $L$, связывание участка отменяется: вершины возвращаются в граф как
    отдельные отношения.
    \item Перестроение графа. Успешно спланированные участки заменяются
    новые вершинами. Строится новый граф соединений
    и пересчитывается $\#\mathrm{csg}$.
\end{enumerate}

После связывания $\#\mathrm{csg}$ пересчитывается.
Если $\#\mathrm{csg} \leqslant \tau$, новый граф передаётся DPsize.
\par

\textit{Шаг 3. Жадный алгоритм.}
Если после свявания $\#\mathrm{csg} > \tau$, оставшийся граф
планируется комбинированным жадным алгоритмом $\mathrm{GOO_{comb}}$.
\par

Если граф имеет несколько компонент связности,
они объединяются тем же $\mathrm{GOO_{comb}}$ в режиме
декартова произведения.
\par

\newpage

\begin{center}
\begin{tikzpicture}[
  >=Latex,
  font=\small,
  every node/.style={align=center},
]
  \tikzset{
    box/.style={draw, rectangle, rounded corners=2pt, minimum width=4cm, minimum height=0.8cm, fill=gray!10},
    pill/.style={draw, rectangle, rounded corners=10pt, minimum width=4.5cm, minimum height=0.7cm, fill=gray!18},
    csg/.style={draw, rectangle, rounded corners=2pt, minimum width=4cm, minimum height=0.8cm, fill=violet!12},
    decide/.style={draw, diamond, aspect=2.5, fill=gray!5, minimum width=3.8cm, inner sep=2pt},
    dpsize/.style={draw, rectangle, rounded corners=2pt, minimum width=3.6cm, minimum height=1cm, fill=teal!18},
    anchorbox/.style={draw, rectangle, rounded corners=2pt, minimum width=4.8cm, minimum height=1cm, fill=orange!22},
    goobox/.style={draw, rectangle, rounded corners=2pt, minimum width=3.6cm, minimum height=1cm, fill=red!15},
    arr/.style={->, thick},
  }

  \node[pill] (start) {\texttt{Список таблиц с условиями соединений}};
  \node[box, below=0.6cm of start] (build) {Построение графа соединений};
  \node[box, below=0.5cm of build] (split) {Разбиение на компоненты связности};

  \node[csg, below=1.2cm of split] (csg1) {Вычисление $\#\mathrm{csg} компоненты$};
  \node[decide, below=0.6cm of csg1] (d1) {$\#\mathrm{csg} \leqslant \tau$};
  \node[dpsize, left=1.8cm of d1] (dp1) {Стандартный DP};

  \node[anchorbox, below=1.2cm of d1] (cont) {Связывание якорных участков\\\scriptsize($\leqslant N$ итераций)};
  \node[decide, below=0.6cm of cont] (d2) {$\#\mathrm{csg} \leqslant \tau$};
  \node[dpsize, left=1.8cm of d2] (dp2) {Стандартный DP};
  \node[goobox, right=1.8cm of d2] (goo) {$\mathrm{GOO_{comb}}$};

  \node[box, below=2.2cm of d2] (cp) {План компоненты};

  \node[draw, dashed, rounded corners=4pt, fit=(csg1) (d1) (dp1) (cont) (d2) (dp2) (goo) (cp), inner sep=0.4cm] (loop) {};
  \node[anchor=north west, font=\footnotesize\itshape, color=gray!50!black, inner sep=4pt] at (loop.north west) {для каждой компоненты связности};

  \node[box, below=0.7cm of loop, minimum width=6cm] (merge) {Слияние компонент: $\mathrm{GOO_{comb}}$ };
  \node[pill, below=0.5cm of merge] (final) {Финальный план};

  \draw[arr] (start) -- (build);
  \draw[arr] (build) -- (split);
  \draw[arr] (split) -- (csg1);
  \draw[arr] (csg1) -- (d1);

  \draw[arr] (d1.west) -- node[above, font=\scriptsize] {да} (dp1.east);
  \draw[arr] (d1.south) -- node[right=1pt, font=\scriptsize] {нет} (cont.north);

  \draw[arr] (cont) -- (d2);

  \draw[arr] (d2.west) -- node[above, font=\scriptsize] {да} (dp2.east);
  \draw[arr] (d2.east) -- node[above, font=\scriptsize] {нет} (goo.west);

  \draw[arr] (dp1.west) -- ++(-0.8,0) |- ([xshift=-1cm, yshift=0.4cm] cp.north) -- ([xshift=-1cm] cp.north);
  \draw[arr] (dp2.south) |- ([yshift=0.4cm] cp.north) -- (cp.north);
  \draw[arr] (goo.south) |- ([xshift=1cm, yshift=0.4cm] cp.north) -- ([xshift=1cm] cp.north);

  \draw[arr] (cp) -- (merge);
  \draw[arr] (merge) -- (final);
\end{tikzpicture}

\smallskip
\captionof{figure}{Схема эвристического поиска порядка соединений: маршрутизация компонент связности по сложности графа $\#\mathrm{csg}$.}\label{fig:heuristic-flow}
\end{center}

\newpage

\begin{center}
\begin{tikzpicture}[
  >=Latex,
  every node/.style={align=center},
]
  \tikzset{
    stepbox/.style={
      draw, rectangle, rounded corners=8pt,
      text width=16cm,
      inner xsep=12pt,
      inner ysep=10pt,
      font=\normalsize\setstretch{1},
      align=center,
    },
    arr/.style={->, thick},
  }

  \node[stepbox, minimum width=16cm, minimum height=3cm] (mark) {%
    Проход по всем вершинам в графе. Для каждой вершины проверяется якорное условие.
    Если выполнено, то вершина с соседями добавляются в новый участок};

  \node[stepbox, below=1.0cm of mark, minimum width=16cm, minimum height=3cm] (filter) {%
    Удаление небезопасных вершин. Если вершина не является якорем, но входит в участок
    и содержит строковое соединение, то она удаляется из участка};

  \node[stepbox, below=1.0cm of filter, minimum width=16cm, minimum height=3cm] (plan) {%
    Каждый участок с помощью DP планируется в новую вершину};

  \node[stepbox, below=1.0cm of plan, minimum width=16cm, minimum height=3cm] (check) {%
    Если кардинальность вершины $> L$, то якорный участок распускается};

  \node[stepbox, below=1.0cm of check, minimum width=16cm, minimum height=3cm] (rebuild) {%
    Граф запроса перестраивается};

  \draw[arr] (mark) -- (filter);
  \draw[arr] (filter) -- (plan);
  \draw[arr] (plan) -- (check);
  \draw[arr] (check) -- (rebuild);
\end{tikzpicture}

\smallskip
\captionof{figure}{Поиск и связывание якорных участков}\label{fig:anchor-contraction}
\end{center}

\newpage

\textbf{Выбор параметров.}
Подход использует четыре параметра: $\tau = 350$, $R_{\max} = 50\,000$,
$\alpha = 100$, $L = 100\,000$.
\par

Порог $\tau = 350$ определяет границу переключения между точным DP
и якорным связыванием.
\par

Значение $R_{\max} = 50\,000$ отсекает все таблицы, сопоставимые по размеру
с таблицами фактов для аналитических схем ($10^5$--$10^7$ строк).
Множитель $\alpha = 100$ требует двухпорядкового различия
между якорем и его крупнейшим соседом, ограничивая якори справочниками
и малыми таблицами, для которых соединения высокоселективны.
Порог связывания $L = 100\,000$ обеспечивает, что новая вершина
не вносит в граф отношение с большой кардинальностью.
Стандартный DP внутри участка выбирает оптимальный план
с высокой вероятностью, поскольку стоимостная модель надёжна
для малых подпланов (см.~секцию <<Оценки в модели стоимости>>).
Проверка $|v \cup N(v)| \leqslant L$ является дополнительной
защитой: если результат связывания оказался неожиданно большим
(признак ошибки оценки), связывание отменяется.
\par

Все четыре значения подобраны экспериментально на JOB.
Дополнительная проверка на вариациях $\tau \in \{250, 300, 350\}$,
$L \in \{25\,000, 50\,000, 100\,000\}$,
$\alpha \in \{50, 100, 200\}$ и $R_{\max} \in \{25\,000, 50\,000, 100\,000\}$
показала, что суммарное отношение $t_{\mathrm{total}}$
к стандартному DPsize остаётся в узком диапазоне
$0{,}86$--$0{,}89$ на JOB и $0{,}11$--$0{,}12$ на JOB-Complex,
то есть подход слабо чувствителен к конкретному выбору в окрестности
выбранных значений.
\par

\textbf{Защита центральных вершин.}
Свойство алгоритма: якорное условие
$|v| < R_{\max} \;\wedge\; |v| \cdot \alpha < \max_{u \in N(v)} |u|$
гарантирует, что центральные вершины (например, в IMDB таблица \texttt{title}
с $|t| = 2{,}5 \cdot 10^6$) \emph{никогда} не становятся якорями.
Центральная вершина может попасть в якорный участок только как сосед якоря~---
но в этом случае она входит в участок \emph{вместе со всеми} соседями якоря,
а не изолируется от части своих соседей.
Это исключает ситуацию, приводящую к деградации у топологической декомпозиции.
\par

Более того, в типичных запросах JOB центральная вершина имеет степень,
превышающую число вершин в якорном участке.
Поскольку участок включает только соседей якоря,
центральная вершина остаётся связанной с незатронутыми вершинами
и сохраняет все параметризованные пути при последующем DPsize.
\par


\textbf{Результаты.}
Суммарные отношения к стандартному DPsize на JOB (113 запросов):
\begin{itemize}
    \item $t_{\mathrm{plan}}$: $0{,}822$ (улучшение $17{,}8\%$),
    \item $t_{\mathrm{exec}}$: $0{,}935$ (улучшение $6{,}5\%$),
    \item $t_{\mathrm{total}}$: $0{,}893$ (улучшение $10{,}7\%$).
\end{itemize}
Из 113 запросов 91 показал улучшение $t_{\mathrm{total}}$ в строгом смысле;
с учётом ожидаемого разброса runtime в $\pm 2\%$ количество запросов,
не уступающих DPsize, составляет 99.
Время планирования улучшилось в 97 запросах, исполнения~--- в 84.
\par

Абсолютные суммарные времена:
\begin{itemize}
    \item DPsize: $t_{\mathrm{plan}} = 52{,}3$ с, $t_{\mathrm{exec}} = 88{,}2$ с,
    $t_{\mathrm{total}} = 140{,}5$ с.
    \item Связывание якорных участков: $t_{\mathrm{plan}} = 42{,}9$ с,
    $t_{\mathrm{exec}} = 82{,}5$ с, $t_{\mathrm{total}} = 125{,}4$ с.
\end{itemize}
Суммарное улучшение~--- $15{,}1$ секунды на 113 запросах,
из которых $9{,}3$ с приходится на $t_{\mathrm{plan}}$
и $5{,}7$ с~--- на $t_{\mathrm{exec}}$.
\par

\textbf{Анализ улучшения $t_{\mathrm{plan}}$.}
Связывание уменьшает число вершин графа соединений,
подаваемого на вход DPsize. Поскольку $\#\mathrm{csg}(G)$ растёт
экспоненциально с числом вершин, удаление даже 2--3 вершин
может сократить $\#\mathrm{csg}$ в несколько раз.
Наибольшая экономия наблюдается на запросах с 13--17 таблицами:
\begin{itemize}
    \item Группа запросов 29a/29b/29c (17 таблиц):
    $t_{\mathrm{plan}}$ снизилось в среднем с $2{,}64$ с
    до $1{,}20$ с (в $2{,}2$ раза). Связывание сократило число вершин
    на 4--5, уменьшив $\#\mathrm{csg}$ ниже порога $\tau = 350$.
    \item Группа 17a--17d (14 таблиц): $t_{\mathrm{plan}}$ снизилось
    с $\sim\!250$ мс до $\sim\!260$ мс на запрос, причём дополнительно
    выиграло $t_{\mathrm{exec}}$ за счёт более удачного порядка.
\end{itemize}
Для запросов с $\leqslant 10$ таблицами $\#\mathrm{csg}$ уже невелик,
и экономия $t_{\mathrm{plan}}$ составляет единицы процентов.
\par

\textbf{Анализ улучшения $t_{\mathrm{exec}}$.}
Улучшение $t_{\mathrm{exec}}$ на $6{,}5\%$ не является целью подхода,
но является его следствием.
Связывание фиксирует высокоселективные соединения (якорь + соседи)
до основного DP.
Внутри участка стандартный DP находит оптимальный план,
поскольку участки малы и оценки кардинальностей для первых соединений надёжны.
Новая вершина, созданная связыванием, несёт точную оценку кардинальности
и фиксирует выбор физического пути
(например, параметризованное индексное сканирование внутри участка).
Это решение \emph{не подвергается пересмотру} на последующих уровнях DPsize,
где ошибки каскадируются и стоимостная модель менее надёжна.
Наиболее яркий пример~--- запрос Q29c: $t_{\mathrm{exec}}$ снизилось
с $3{,}35$ с до $84$ мс (в $40$ раз), поскольку якорный план избежал
четырёхуровневого вложенного сканирования с промежуточным результатом в $268$ тысяч строк,
который DPsize выбрал из-за заниженной оценки кардинальности
на средних уровнях при построении плана.
\par

\textbf{Анализ деградаций.}
Из 22 запросов JOB, на которых $t_{\mathrm{total}}$ ухудшилось,
20 имеют деградацию не более $12\%$, и суммарные потери составляют
$2{,}9$ секунды (против $18{,}0$ секунды суммарного выигрыша
на остальных 91 запросе).
Наибольшие деградации сосредоточены в группе запросов 31a/31b/31c
(14 таблиц): $t_{\mathrm{total}}$ ухудшилось на $11\%$--$43\%$
($+387$, $+503$, $+993$ мс соответственно).
Причина: связывание выделило участок, в котором локальный DPsize
по-другому распределил физические операторы,
что привело к замене Hash хэш соединения на менее эффективное вложенное сканирование на верхнем
уровне.
\par

Среди 16 запросов с деградацией $t_{\mathrm{plan}}$
причиной является накладной расход на построение графа соединений,
вычисление $\#\mathrm{csg}$ и планирование участков, который не окупается
на запросах с $\#\mathrm{csg} \leqslant \tau$, где стандартный DPsize
и так быстр (типично $5$--$25$ мс).
\par

\textbf{Результаты на JOB-Complex.}
На 30 запросах JOB-Complex суммарные отношения к стандартному DPsize:
\begin{itemize}
    \item $t_{\mathrm{plan}}$: $0{,}867$ (улучшение $13{,}3\%$),
    \item $t_{\mathrm{exec}}$: $0{,}073$ (улучшение $92{,}7\%$),
    \item $t_{\mathrm{total}}$: $0{,}117$ (улучшение $88{,}3\%$).
\end{itemize}
Абсолютные времена: DPsize $t_{\mathrm{total}} = 648{,}8$~с
(из них $612{,}7$ с на исполнение), якорный подход $t_{\mathrm{total}} = 75{,}9$~с
(из них $44{,}5$ с на исполнение).
Суммарная экономия~--- $572{,}9$ секунды на 30 запросах,
причём $568{,}2$ с приходится на $t_{\mathrm{exec}}$.
В строгое улучшение $t_{\mathrm{total}}$ показали 22 запроса из 30.
\par

Доминирующий вклад в выигрыш даёт запрос Q12:
DPsize строит план с верхним вложенным сканированием по строковому предикату
\texttt{chn.name\_pcode\_nf = ak.name\_pcode\_nf},
что приводит к промежуточному результату в $\sim\!6{,}5 \cdot 10^8$ строк
и времени исполнения $546{,}9$~с.
Якорный подход связывает участок вокруг \texttt{company\_name}
и заменяет вложенное соединение на хэш соединение по тому же предикату на ранней
стадии плана, что снижает $t_{\mathrm{exec}}$ до $668$ мс (в $\sim\!820$ раз).
Аналогичный механизм избегания вложенного сканирования по
строковому предикату срабатывает на запросах Q5, Q11, Q17, Q20:
$t_{\mathrm{exec}}$ снижается на 1--2 порядка.
\par

Наибольшая оставшаяся деградация~--- Q21
($t_{\mathrm{total}}$: $4{,}79$~с $\to 6{,}57$~с, $+37\%$): связывание
выделило участок, для которого локальный DPsize выбрал план,
проигрывающий на верхнем уровне основного DP.
\par

\textbf{Результаты на истинных кардинальностях.}
В ходе исследования были также получены данные в условиях истинных кардинальностей. Для каждого запроса перебирались все связанные подграфы,
каждому такому подграфу сопоставлялся подзапрос,  для него вычислялась кардинальность. Таким образом каждому запросу из JOB и
JOB-Complex, соответстует множество кардинальностей для всех его подзапросов. Было разработано расширение, которое подставляет истинные
значения при оценке кардинальностей в планировщик PostgreSQL. Для эффективного доступа используется хэш таблица.
\par

На наборе JOB при истинных кардинальностях суммарное
время $t_{\mathrm{plan}}$ сократилось с $47{,}5$~с до $43{,}1$~с
($-9{,}3\%$), $t_{\mathrm{exec}}$ практически не изменилось
($25{,}4$~с против $26{,}0$~с, $+2{,}2\%$),
а $t_{\mathrm{total}}$~--- с $72{,}9$~с до $69{,}0$~с ($-5{,}3\%$).
Якорный подход не уступает стандартному DPsize (с учётом разброса $\pm 5\%$)
на 96 запросах из 113. Оставшийся выигрыш сосредоточен в планировании:
на больших графах (12--17 таблиц) связывание уменьшает $\#\mathrm{csg}$
и сокращает фазу DP.
На наборе JOB-Complex отношения составили
$0{,}892$ по $t_{\mathrm{plan}}$, $0{,}574$ по $t_{\mathrm{exec}}$
и $0{,}795$ по $t_{\mathrm{total}}$
(абсолютная экономия $9{,}8$~с из общего $t_{\mathrm{total}} = 47{,}8$~с).
\par

При истинных кардинальностях главная компонента выигрыша смещается:
относительная экономия $t_{\mathrm{plan}}$ становится меньше, чем при
статистических оценках ($9{,}3\%$ против $17{,}8\%$ на JOB), поскольку
истинные значения уменьшают абсолютное время DPsize.
$t_{\mathrm{exec}}$ при этом перестаёт быть источником выигрыша на JOB:
если стоимостная модель получает точные оценки, она находит хорошие порядки, 
и связывание уже не приносит дополнительной
выгоды на уровне исполнения.
\par

Иллюстративный пример~--- запрос Q29b из JOB: на истинных кардинальностях
DPsize строит план за $2{,}56$~с при $t_{\mathrm{exec}} = 6$~мс;
якорный подход~--- за $1{,}21$~с при $t_{\mathrm{exec}} = 35$~мс.
Граф запроса содержит около $17$ вершин с несколькими малыми
справочниками-якорями (\texttt{role\_type}, \texttt{info\_type},
\texttt{comp\_cast\_type}, \texttt{keyword}); их соседи
образуют непересекающиеся участки. После связывания
$\#\mathrm{csg}$ нового графа падает ниже порога~$\tau$, и DPsize
завершает перебор за время, сравнимое с накладными расходами на сам
каскад. Выигрыш по $t_{\mathrm{total}}$~--- $1{,}32$~с ($-51\%$).
\par

Наибольшая оставшаяся деградация на JOB-Complex при истинных
кардинальностях~--- Q28 ($t_{\mathrm{total}}$: $598 \to 1059$ мс, $+77\%$):
связывание создало виртуальную вершину, для которой основной DPsize
выбрал вложенное сканирование вместо хэш соединения,
доступного в исходном графе.
\par

